\documentclass[exercice]{thedocument}%
\startdatakeys
\source{}
\BO{2026}
\cycle{4}
\competencesDuSocle{RA,MO}
\objectifsApprentissage{C4D1bg,C4D2ab,C4D1bf}
\automatismes{C4A1cf,C4A3bc}
\enddatakeys
\begin{document}%
Le dimanche, Lucas achète 3 beignets à la boulangerie, pour le prix de 7 euros. Le
lendemain, il achète 9 beignets pour 20 euros.
\begin{enumerate}
    \item Quelles sont les deux grandeurs citées dans cette énoncé ?
    \item Sont-elles proportionnelles ? Justifier la réponse.
\end{enumerate}
\end{document}
